\section{Public Dataset Selection}
Dataset selection was governed by four requirements for the primary association analyses: (i) survey-weighted population representativeness; (ii) objective pure-tone audiometry; (iii) tinnitus and/or hearing-difficulty items; and (iv) covariates for noise exposure, psychosocial stress, and general health. Only KNHANES and NHANES jointly satisfy all four at population scale. Supplementary Table~S1 summarizes the selection, and Supplementary Table~S2 maps target variables to each source.

\subsection{Primary Development Dataset: KNHANES}
KNHANES is selected as the primary development dataset because it is nationally representative for Korea and has previously supported tinnitus and hearing-loss analyses \citep{park2014tinnitus,joo2015hrqol}. Its advantages include Korean population relevance, complex-survey weights, pure-tone audiometry in applicable cycles (otologic examination cycles 2010--2012), tinnitus and hearing-difficulty questionnaire items in selected cycles, health and behavioral covariates, and psychosocial or stress-related self-report variables (e.g., perceived-stress and depressed-mood items). It also provides a realistic bridge to future Korean hospital cohorts, including the planned prospective clinical extension. Because item wording and audiometry availability differ across cycles, the analysis prespecifies a per-cycle codebook review and restricts pooling to cycles with comparable definitions.

The limitations are equally important. KNHANES is not a workplace cohort, and the audiometric analyses are cross-sectional. It can therefore evaluate \emph{associations} among perceived stress, work-related variables, tinnitus, and audiometric outcomes, but cannot by itself establish that occupational stress \emph{caused} tinnitus or hearing loss. Treatment timing and recovery after sudden hearing loss are not captured.

\subsection{External-Validation Dataset: NHANES}
NHANES is selected as the primary external-validation dataset because it provides openly downloadable U.S. public-use files, pure-tone audiometry and tympanometry in selected cycles (2011--2012, 2015--2016, 2017--2018), occupational and recreational noise-exposure items, tinnitus and hearing questionnaires in selected cycles, and extensive demographic and health covariates. NHANES enables testing whether feature definitions and risk models derived from KNHANES retain discrimination and calibration across population, language, health-system, and occupational contexts. As with KNHANES, complex-survey weights, strata, and primary sampling units must be honored, and cycle-specific age eligibility for the audiometry module must be respected.

\subsection{Prospective Clinical Extension (Hospital-Based Cohort)}
A prospective otolaryngology/audiology cohort is specified to address what public cross-sectional data cannot---treatment timing, SSNHL recovery, longitudinal tinnitus burden, and rehabilitation---and to supply deidentified audiograms and clinical narratives for the companion clinical validation and safety evaluations (Supplementary Methods~S2). Raw clinical data are not redistributed; only derived, deidentified, or synthetic artifacts are shared, subject to IRB approval.
